\documentclass{article}
\usepackage{graphicx} % Required for inserting images
\usepackage{amsmath}
\usepackage[english, russian] {babel}
\usepackage[utf8]{inputenc}
\usepackage[T2A]{fontenc}
\usepackage{minted}
\usepackage{float}
\usepackage{amssymb}

\title{Финансовая грамотность}
\author{silvia.lesnaia}


\begin{document}

\maketitle

    \large{21.02.26}

    \vspace{2mm}
    \Large{\fbox{Лекция}}
    \section{Деньги сущность}
    \subsection{Сбережения и потребления}

    Сбережения - приумножения текущих доходов, с расчетом на удовлетворения 
    потребностей в будущем. Пропорции сбережний и потребления определяют зависимости 
    от текущего дохода.

    Перспектива доходов в будущем

    Предпочтения сегодняшенго потребления по сравнению с будущем

    \subsection{Функция денег}

    \begin{itemize}
        \item Мера стоимости  - деньги аналогичны едециницам зимереня и позволяет овценивать 
        различные блага
        \item Средства обращения - используется при переходе товаров из рук в руки, 
        позволяя избежать бартера 
        \item Средства платежа - оплата услуг, по сути то что нельзя оплатить бартером
        \item Средства накопления 
        \item Мировые деньги
    \end{itemize}
    
    СДР - 

    \subsection{Рынок денег}

    Предоставлется улсгуи по продажи днежных средсвт, движение. 

    \vspace{1cm}

    \Large{\fbox{Практика}}

    Ну сегодня мы просто беседуем

    \vspace{1cm}
    \large{07.03.26}

    \vspace{2mm}
    \Large{\fbox{Лекция}}

    \subsection{Рынок валют} 

    Если деньги сущетсвуют деньги и не включаются в мир, то они просто 
    локальные деньги, а если включаются то становят ся валютой. 

    Пример коллективной валютой: евро

    Валюта классифицируется по конветритруемсямая: 

    \begin{itemize}
        \item ковертируемость
        \item не конвертирумесь
        \item частино конвертируемой
    \end{itemize}


\end{document}